\section{Utterance and articulating}

Owner ruling 2026-09-21, ``the admission boundary, restored'': the utterance's
two hands are \emph{types}, not positions read off a bare pair of marks
afterward. A token is a mark carried on the left hand; a claim is a mark
carried on the right hand; the utterance is the pair of the two typed
carriers, and it is the theory's sole admission boundary --- the one place a
mark is let in.

\begin{definition}[Token]
  \label{def:token}
  \lean{Obsmetrologia.Token}
  \leanok
  \uses{def:mark}
  A \emph{token} is a mark carried on the left hand: the one-field structure
  $\{\text{mark} : \texttt{Mark}\}$. It is the identity-carrying side of an
  utterance --- \emph{which one} --- distinguished from a claim by its own
  type, not recovered afterward by position.
\end{definition}

\begin{definition}[Claim]
  \label{def:claim}
  \lean{Obsmetrologia.Claim}
  \leanok
  \uses{def:mark}
  A \emph{claim} is a mark carried on the right hand: the one-field structure
  $\{\text{mark} : \texttt{Mark}\}$, dual to \cref{def:token}. It is the
  content-carrying side of an utterance --- \emph{what is said}.
\end{definition}

\begin{definition}[Utterance]
  \label{def:utterance}
  \lean{Obsmetrologia.Utterance}
  \leanok
  \uses{def:token, def:claim, thm:kind}
  An \emph{utterance} is the ordered pair of a token and a claim --- the two
  typed hands, token first, claim second --- and the domain of articulating.
  The ordering of the pair is the kind distinction of \cref{thm:kind}, now
  carried at the level of the two hands' own types.
\end{definition}

\begin{definition}[Articulating]
  \label{def:articulating}
  \lean{Obsmetrologia.Articulating}
  \leanok
  \uses{def:utterance}
  \emph{Articulating} is the base morphism: the first arrow of the generative
  unfolding, from which every other morphism is reached. Inductively it is
  either an innermost utterance (\texttt{utter}, the recursion's primitive
  floor) or an utterance articulated over a further articulating
  (\texttt{pair}, the pair $[\text{utterance}, \text{articulating}]$). This is
  the second and last admission of free data in Obsmetrologia.
\end{definition}

\begin{proposition}[Floor]
  \label{thm:floor}
  \lean{Obsmetrologia.Articulating.floor}
  \leanok
  \uses{def:articulating}
  Every articulating bottoms out at a primitive utterance, its \emph{floor},
  recovered by descending through the \texttt{pair} chain. Totality of
  \texttt{floor} witnesses that grounding in primitive marks yields finite
  inhabitants.
\end{proposition}

\begin{lemma}[Articulating is one-way]
  \label{lem:articulating-oneway}
  \lean{Obsmetrologia.Articulating.utter_ne_pair}
  \leanok
  \uses{def:articulating}
  A primitive floor is never a nested pair:
  \texttt{utter}~$u \neq$ \texttt{pair}~$v~a$. No arrow collapses the one-way
  structure; to descend is to derive, never to un-derive. This is the
  Curry--Howard directionality the filepath-as-proof discipline rests on.
\end{lemma}
